% !TeX root = ../../../金融经济学_王江_讲义及习题.tex
% 借鉴 431 计算讲义的两资产曲线与三资产两两组合画法。
% 中图取三只证券均值为 1、2、3，协方差矩阵为单位阵；左图只取其中 A、C。
% 横向绘图坐标为 sqrt(3)*sigma，纵向绘图坐标为均值减 0.45。
\begin{tikzpicture}[>=Stealth,x=1.06cm,y=.9cm,font=\small]
 \foreach \panel/\offset/\heading in {
   1/0/{(a) 两只证券 $N=2$},
   2/4.05/{(b) 三只证券 $N=3$},
   3/8.1/{(c) 一般情形 $N\ge3$}} {
  \begin{scope}[xshift=\offset cm]
   \def\minrisk{1}
   \ifnum\panel=1
    \def\minrisk{1.5}
   \fi
   \ifnum\panel>1
    \fill[blue!8] plot[domain=-1.25:1.25,samples=65,variable=\t]
     ({sqrt(1+1.5*\t*\t)},{1.55+\t})--(2.9,2.8)--(2.9,.3)--cycle;
   \fi
   \draw[->] (0,0)--(3.05,0) node[right] {$\sigma_p$};
   \draw[->] (0,0)--(0,3.02) node[above] {$\bar r_p$};
   \draw[gray!75,thick,domain=-1.25:0,samples=35,variable=\t]
    plot ({sqrt(\minrisk+1.5*\t*\t)},{1.55+\t});
   \draw[red!70!black,very thick,domain=0:1.25,samples=35,variable=\t]
    plot ({sqrt(\minrisk+1.5*\t*\t)},{1.55+\t});
   \node at (1.55,-.57) {\heading};
   \fill ({sqrt(\minrisk)},1.55) circle(1.4pt) node[left] {\scriptsize MVP};
  \end{scope}
 }
 \foreach \t/\asset in {-1/A,1/C} {
  \fill[blue!65!black] ({sqrt(3)},{1.55+\t})
   circle(1.4pt) node[right] {$\asset$};
 }
 \begin{scope}[xshift=4.05cm]
  \foreach \low/\high in {1/2,1/3,2/3} {
   \draw[blue!70!black,dashed,domain=0:1,samples=32,variable=\t]
    plot ({sqrt(3*((1-\t)^2+\t^2))},{\low+(\high-\low)*\t-.45});
  }
  \foreach \mean/\asset in {1/A,2/B,3/C} {
   \fill[blue!65!black] ({sqrt(3)},{\mean-.45})
    circle(1.4pt) node[right] {$\asset$};
  }
 \end{scope}
 \begin{scope}[xshift=8.1cm]
  \coordinate (frontier) at ({sqrt(1+1.5*.65^2)},2.2);
  \draw[->,blue!65!black,thick] (frontier)--(2.78,2.2)
   node[midway,above] {\scriptsize 固定均值};
  \fill[red!70!black] (frontier) circle(1.4pt) node[above left] {$P_p$};
  \node[blue!65!black] at (2.12,1.35) {\scriptsize 可行区域};
 \end{scope}
\end{tikzpicture}
